% ch3-facts.tex
\chapter{Assumed, Verified, Inferred, Checked: Where Facts Come From}
\label{ch:facts}

\begin{answer}
ulm.lam proves nothing about the entries of a matrix. Facts enter in two
places only, the parameter types and \code{assume}, and are trusted there.
Inference rules derive new facts from them; these are theorems about real
matrices, valid if the declarations are. Verifiers check that the IR is
consistent with the rules, never with the data. At run time the Julia (or
Python) runtime checks the array sizes, and each factorization checks the
one fact it depends on: \code{potrf} detects a matrix that is not positive
definite in floating point, \code{getrf} an exactly singular one. Symmetry
of an input, full column rank and conditioning are not checked.
\end{answer}

% ---------------------------------------------------------------------------
\section{The facts and how they imply each other}
\label{sec:facts-lattice}

A property set is a conjunction of atoms. Its meaning is fixed by one C++
function, \code{closure} (\srcref{lib/UlmLamDialect.cpp}{97}{123}), which
adds everything that the given atoms imply (\cref{fig:implications}).

\begin{figure}[h]
\centering
\begin{tikzpicture}[font=\sffamily\small,
  p/.style={draw=accent,fill=accent!6,rounded corners=2pt,inner sep=3pt,minimum height=0.6cm},
  c/.style={draw=accent,circle,fill=white,inner sep=1.2pt,font=\sffamily\scriptsize},
  arr/.style={-{Stealth[length=1.8mm]},draw=accent!80!black,semithick}]
  \node[p] (id) at (0,0) {identity};
  \node[p] (spd) at (-4.0,-1.4) {spd};
  \node[p] (perm) at (0.6,-1.4) {permutation};
  \node[p] (tri) at (3.6,-1.4) {lower\_tri, upper\_tri};
  \node[p] (ud) at (3.9,-2.8) {unit\_diag};
  \node[p] (spsd) at (-5.4,-2.8) {spsd};
  \node[p] (ns) at (-2.6,-2.8) {non\_singular};
  \node[p] (orth) at (0.6,-2.8) {orthogonal};
  \node[p] (sym) at (-5.4,-4.2) {symmetric};
  \node[p] (fr) at (-0.4,-4.2) {full\_rank};
  \node[p] (sq) at (-3.8,-5.6) {square};
  \node[c] (and) at (-1.4,-5.6) {\(\wedge\)};
  \draw[arr] (id) -- (spd); \draw[arr] (id) -- (perm); \draw[arr] (id) -- (tri);
  \draw[arr] (id) -- (ud);
  \draw[arr] (spd) -- (spsd); \draw[arr] (spd) -- (ns);
  \draw[arr] (spsd) -- (sym); \draw[arr] (sym) -- (sq);
  \draw[arr] (ns) -- (sq); \draw[arr] (ns) -- (fr);
  \draw[arr] (perm) -- (orth); \draw[arr] (orth) -- (fr);
  \draw[arr] (sq) -- (and); \draw[arr] (fr) -- (and);
  \draw[arr] (and) -- (ns);
\end{tikzpicture}
\caption{The implications applied by \code{closure}. The circle is a
conjunction: square and full rank imply non-singular. A static square shape
adds \code{square}. \code{lower_tri} together with \code{upper_tri} is
printed as \code{diagonal}.}
\label{fig:implications}
\end{figure}

Three conventions make this a sound basis for reasoning:
\begin{itemize}
\item \textbf{Absence means unknown.} A set without \code{spd} says nothing
  about definiteness; it does not say ``indefinite''. The empty set
  (printed as no encoding) means ``some matrix''.
\item \textbf{Dropping a fact is always sound.} A value with
  \code{\{spd\}} may be used where \code{\{symmetric\}} or \code{\{\}} is
  expected. Rewrites use this to keep result types unchanged when they find
  a better-typed replacement (they insert an \code{assume} that only drops
  facts, \code{replaceKeepingType},
  \srcref{lib/UlmLamDialect.cpp}{567}{576}).
\item \textbf{Two bits are not facts about the value.} \code{lower_stored}
  and \code{upper_stored} say that only one triangle of the \emph{memory}
  is meaningful. They are never inferred and never propagated to a result,
  and dropping them would be wrong; the comment in
  \srcref{include/ulm_lam/UlmLamAttrs.td}{10}{34} spells this out.
\end{itemize}

\begin{theory}[ (\FGs{5.1}, \FGs{5.2})]
The property sets, ordered by implication, form a finite lattice; ``absence
means unknown'' makes each set an over-approximation of the matrices it
describes, which is the abstraction of~\cite{cousot} and
\FG{5.2}{Abstract interpretation: Galois connections, widening, narrowing}.
\code{closure} plays the role of a reduction (it adds what the other atoms
imply, without changing the set of matrices described), and each op's rule
is a sound abstract transfer function. Because a function without loops is
a DAG (directed acyclic graph) of SSA values, one pass in program order
computes the result; no fixpoint iteration
(\FG{5.1}{Lattices, monotone frameworks and fixpoints (a sharper recap)}) is
needed. Loops are the exception, discussed next.
\end{theory}

\paragraph{Loops: facts are re-declared, not inferred} A loop-carried
matrix enters the loop body as a block argument. The Julia frontend gives
every loop-carried value a type without properties (\code{plainvalue} in
\code{loop}, \srcref{julia/UlmLam/src/frontend.jl}{408}{434}), and
\code{ulm-lam-infer-properties} refines op results only, not block
arguments. A fact that holds in every iteration therefore has to be
re-declared, as the Kalman filter in the README does:
\begin{lstlisting}[style=jl]
loop(steps, x0, P0) do i, x, P
    P = assume(P, :symmetric)
    ...
\end{lstlisting}
Proving that \(P\) stays symmetric would need a fixpoint over the loop: start
from the facts of \(P_0\), apply the body's transfer functions, and
intersect with the facts at the loop head until nothing changes. That is
exactly the iterative dataflow of \FGs{5.1}; ulm.lam does not implement it
and asks the programmer instead.

% ---------------------------------------------------------------------------
\section{Declared facts: parameters and \texttt{assume}}
\label{sec:facts-declared}

\begin{steps}
\item[Input] \code{assume(x, :spd)} in Julia (\code{lam.assume(x, "spd")} in Python).
\item[Representation] \code{\%y = ulm_lam.assume \%x : tensor<...> -> tensor<..., #ulm_lam.props<spd>>}.
\item[Rule] none: the result has the properties of its type. The verifier only checks that shape and element type are unchanged.
\item[Output] a value that later rules treat as SPD.
\item[Theory] an unchecked cast into a refinement type; the trusted boundary of \FG{1.7}{Refinement types}.
\end{steps}

The ODS description (\srcref{include/ulm_lam/UlmLamOps.td}{135}{154}) says:
``Properties the operand does not have are \emph{assertions}: with
\code{checked} they are verified at run time [\ldots], otherwise trusted.
This is the only op that may add facts that inference cannot derive.'' At
this commit the \code{checked} variant is not implemented
(\cref{sec:infra-analogy}); the Julia docstring of \code{assume}
(\srcref{julia/UlmLam/src/frontend.jl}{351}{357}) states the actual
behavior: ``Trusted; a factorisation that finds it wrong returns
info \(>0\) (\code{InfoError}).''

\paragraph{What a declaration changes} The same GLS formula with different
declarations (\code{outputs/variants.out}, shortened):

\begin{textout}
gls_nofr(X::n×m, M::n×n spd, y::n)                 # X without :full_rank
  potrf    L1 ← chol(M)
  copy     T1 ← X
  trsm     T1 ← L1⁻¹ T1
  syrk     T2 ← T1ᵀ T1                       one triangle only
  symm.    T2 ← sym(T2)                      copy one triangle to the other
  getrf    LU1 ← lu(T2)                      P T2 = L U, LU
  ...
gls_assume(X::n×m, M::n×n spd, y::n)               # S = assume(X' * inv(M) * X, :spd)
  potrf    L1 ← chol(M)
  copy     T1 ← X
  trsm     T1 ← L1⁻¹ T1
  syrk     T2 ← T1ᵀ T1                       one triangle only
  potrf    L2 ← chol(T2)                     T2 = L2 L2ᵀ, Cholesky
  ...
\end{textout}

Without \code{full_rank}, \(S = X\T M\inv X\) can only be inferred SPSD, so
\(S\inv\) is not mapped to Cholesky: kernel selection completes the
symmetric matrix (\code{sycomplete}) and factorizes it by LU. With the
assumption \(S\) is SPD again and the Cholesky path returns. All variants
compute the same \(\beta\) for the example data; they differ in cost and in
what they check.

\paragraph{A false declaration} \(M\) declared SPD but not symmetric
(\(M_{21}=100\) in Julia's lower triangle, \code{outputs/variants.out}):

\begin{textout}
ulm.lam: [1.2739726027397262, 0.9452054794520546]
Julia with Symmetric(Mns, :U): [1.2739726027397262, 0.9452054794520547]
Julia with Mns as given:       [1.2726368159203976, 0.9552238805970151]
\end{textout}

\code{potrf} reads one triangle (in Julia's terms the upper one: the IR
holds the transpose) and never looks at the other, so ulm.lam computes the
GLS estimate for the symmetric matrix that the upper triangle defines. No
error is reported; the result is the answer to a different problem. The
origin of symmetry is \Odecl{} and its check is \Onone.

\begin{limit}
An assumption is not a run-time check. \code{assume(A, :spd)} and the
parameter type \code{Mat(:n, :n, :spd)} make the compiler \emph{act} on SPD;
they do not test it. A Cholesky factorization of \(A\) does test positive
definiteness of the stored triangle, as a side effect of computing the
factor. Whether such a factorization happens depends on the formula: a
matrix that is declared SPD but only multiplied (\(A x\)) is never tested.
\end{limit}

% ---------------------------------------------------------------------------
\section{Inferred facts: the rules}
\label{sec:facts-inferred}

\Cref{tab:rules} lists every inference rule, read from
\srcref{lib/UlmLamDialect.cpp}{192}{357} and
\srcref{lib/UlmLamDialect.cpp}{900}{939}. All are facts of linear algebra
over \(\R\); each is applied to the \emph{closures} of the operand sets.

\begin{table}[h]
\centering\small
\begin{tabularx}{\linewidth}{@{}P{3.0cm}L@{}}
\toprule
\textbf{Op} & \textbf{Result facts} \\
\midrule
\code{mul} \(F_1\cdots F_k\) & the facts shared by all factors among
  \code{lower_tri}, \code{upper_tri}, \code{unit_diag}, \code{orthogonal},
  \code{permutation}, \code{non_singular}, \code{identity}; \code{zero} if
  any factor is zero; \code{full_rank} for \(NA\) or \(AN\) with \(N\)
  non-singular and \(A\) full rank (two factors); and the palindrome rule:
  \(Z\T Z\) SPSD (SPD if \(Z\) has full column rank),
  \(Z\T M Z\) symmetric / SPSD / SPD if \(M\) is \\
\code{add} \(A_1+\cdots+A_k\) & the facts shared by all terms among
  \code{symmetric}, \code{spsd}, \code{lower_tri}, \code{upper_tri}; SPD if
  all are SPSD and one is SPD \\
\code{scale} \(\alpha A\) & triangularity, symmetry, zero; SPD/SPSD if
  \(\alpha\) is a \emph{constant} \(>0\); rank facts if \(\alpha\ne 0\);
  zero if \(\alpha = 0\) \\
\code{transpose} & everything, with \code{lower_tri} and \code{upper_tri}
  exchanged; storage bits dropped \\
\code{inverse} & triangularity, unit diagonal, symmetric, SPD, orthogonal,
  permutation, identity, non-singular, full rank, square \\
\code{identity} & \code{identity} \\
\code{assume} & the properties of its result type (declared) \\
\midrule
\code{ulm_lamk.trsm} & \code{full_rank} if the right-hand side has full rank \\
\code{ulm_lamk.syrk} & \code{spsd}; \code{ulm_lamk.gemmtr}: \code{symmetric} \\
other kernels & nothing (\code{potrf}'s result holds a factor, not the matrix) \\
\bottomrule
\end{tabularx}
\caption{The inference rules of ulm.lam at commit \texttt{\shortcommit}.}
\label{tab:rules}
\end{table}

Two of the rules show how much depends on details:
\begin{itemize}
\item \code{scale} keeps SPD only for a positive \emph{constant}. For a
  scalar known only at run time (\code{Float} parameter), the sign is
  unknown, so \(\sigma^2 K\) loses SPD. This is why the README writes
  \code{assume(K + s2 * I, :spd)} for a Gaussian process.
\item The palindrome rule compares operands by \emph{identity}, not by
  value: \(X\T M\inv X\) is recognized because both outer factors are the
  same SSA value, once transposed. Two different values that happen to hold
  equal matrices are not recognized. The interface description says so
  (``Operand \emph{identity} matters, not only operand properties'',
  \srcref{include/ulm_lam/UlmLamInterfaces.td}{11}{27}).
\end{itemize}

\paragraph{Why inferred facts matter} They choose the kernels.
\Cref{tab:variants} summarizes \code{outputs/variants.out}.

\begin{table}[h]
\centering\small
\begin{tabularx}{\linewidth}{@{}P{3.6cm}P{3.0cm}L@{}}
\toprule
\textbf{Declarations} & \textbf{Inferred for \(S\)} & \textbf{Kernels for \(M\inv\) and \(S\inv\)} \\
\midrule
\(X\) full rank, \(M\) SPD & SPD & \code{potrf}, \code{potrf}; \code{syrk} for \(S\) \\
\(X\) general, \(M\) SPD & SPSD & \code{potrf}; \code{syrk}, \code{sycomplete}, \code{getrf} \\
\(X\) full rank, \(M\) symmetric & symmetric & \code{getrf}, \code{getrf}; \code{gemmtr} for \(S\) \\
\(X\) full rank, \(M\) general & none & \code{getrf}, \code{getrf}; \code{gemm} for \(S\) \\
\(X\) general, \(M\) SPD, \code{assume(S, :spd)} & SPD (declared) & \code{potrf}, \code{potrf}; \code{syrk} \\
\bottomrule
\end{tabularx}
\caption{The same formula under different declarations.}
\label{tab:variants}
\end{table}

\begin{limit}
Property inference is not a general theorem prover. It applies a fixed set
of rules, each local to one op, to a fixed vocabulary of fifteen atoms.
It cannot use facts about sizes (\(n \ge m\)), cannot relate two different
values, and cannot prove anything that needs algebra across ops, for example
that \(A + A\T\) is symmetric (\code{add} only keeps facts that all terms
share). Where a rule does not apply, the fact is unknown, and the program is
still compiled, just with more general kernels.
\end{limit}

% ---------------------------------------------------------------------------
\section{Facts established by analyses}
\label{sec:facts-analysis}

Some facts are not about matrices but about the program. They are found by
analyses that inspect the IR:

\begin{itemize}
\item \textbf{Which size a dynamic dimension has.} \code{dimOf}
  (\srcref{lib/UlmLamDialect.cpp}{1269}{1357}) follows a value back through
  ulm.lam ops, DPS destinations, loops, branches and calls to the function
  parameter that \code{ulm_lam.dims} names. \Cref{sec:shapes-dims}.
\item \textbf{Whether an earlier product may be reused.} Kernel selection
  reuses \(W\) only if its definition properly dominates the use
  (\code{dom.properlyDominates}, \srcref{lib/UlmLamPasses.cpp}{1080}{1091}).
\item \textbf{Whether an op is loop invariant.} LICM moves an op out of a
  loop if its operands are defined outside and it is pure
  (\cref{sec:effects-licm}).
\item \textbf{Whether a write may be in place.} One-Shot Bufferize's
  read-after-write conflict analysis (\cref{ch:storage}).
\end{itemize}

These are \Oana{} facts: derived from the program text, sound for all
inputs, and independent of the matrix entries.

% ---------------------------------------------------------------------------
\section{Facts checked at run time}
\label{sec:facts-runtime}

\begin{itemize}
\item \textbf{Array sizes}, in the Julia runtime before the C call
  (\code{bind_sizes!}, \srcref{julia/UlmLam/src/runtime.jl}{326}{339}):
  every run-time size symbol must have one value across all arguments, and
  every fixed size must match. The generated C function takes the sizes as
  arguments and does not check them.
\item \textbf{Positive definiteness}, by \code{potrf}: LAPACK reports
  \(\mathrm{info}=i>0\) if the leading minor of order \(i\) is not
  positive definite \emph{in floating point}. The compiled function returns
  the number of the check (\cref{sec:effects-info}).
\item \textbf{Exact singularity}, by \code{getrf}: \(\mathrm{info}=i>0\)
  if \(U_{ii}\) is exactly zero. A nearly singular matrix passes.
\item \textbf{Workspace size}, if the caller provides the workspace:
  \(\mathrm{info} = -i\) when it is too small.
\end{itemize}

What is not checked anywhere: that a matrix declared symmetric (or SPD) is
symmetric; that \(X\) has full column rank or that \(n\ge m\); that a
declared triangular matrix has zeros in its other triangle (the kernels
never read it); the conditioning of anything.

% ---------------------------------------------------------------------------
\section{Three ways to know that a matrix is SPD}
\label{sec:facts-three}

The task statement of this companion contrasts three situations. With the
GLS example they are:

\begin{center}\small
\begin{tabularx}{\linewidth}{@{}P{3.2cm}P{3.6cm}L@{}}
\toprule
\textbf{Situation} & \textbf{Origin} & \textbf{Consequence if the fact is false} \\
\midrule
\code{M::Mat(:n, :n, :spd)} & \Odecl{}, trusted & \code{potrf} of \(M\) fails (\code{info 1}) if \(M\) is not positive definite; wrong \(\beta\), no error, if \(M\) is not symmetric \\
\(S = X\T M\inv X\) is SPD & \Oinf{} from \(M\) SPD and \(X\) full column rank & cannot be false over \(\R\) if the premises hold; in floating point \code{potrf} of \(S\) may still fail (\cref{sec:kernels-fp}) \\
\code{potrf} returns \(\mathrm{info}>0\) & \Orun{} & the call returns \code{info}; Julia throws \code{InfoError}; outputs unspecified \\
\bottomrule
\end{tabularx}
\end{center}

The first is an assumption, the second a conditional theorem, the third a
measurement. Only the third looks at numbers, and only at the numbers that
the chosen kernels happen to factorize.
